\documentclass{article}

% camera-ready version
\usepackage[final]{neurips_2024}

\usepackage[utf8]{inputenc}
\usepackage[T1]{fontenc}
\usepackage{hyperref}
\usepackage{url}
\usepackage{booktabs}
\usepackage{amsfonts}
\usepackage{amssymb}
\usepackage{nicefrac}
\usepackage{microtype}
\usepackage{xcolor}
\usepackage{graphicx}
\usepackage{amsmath}
\usepackage{multirow}
\usepackage{array}

\title{Disentangling Feature Absorption: Confound Resolution, Cross-Domain Probing, and Scaling Surfaces for SAE Quality Assessment}

\author{%
  Anonymous Author(s)
}

\begin{document}

\maketitle

\begin{abstract}
Feature absorption---where an SAE latent silently fails to fire because a more specific, co-occurring latent subsumes its role under the sparsity objective---is the most documented failure mode of sparse autoencoders (SAEs) for mechanistic interpretability, yet its causal relationship to downstream quality has never been tested with proper confound control.
We apply epidemiological causal inference methods (partial correlation, Baron-Kenny mediation, Rosenbaum sensitivity analysis) to 48 Gemma Scope SAEs and find that controlling for $L_0$ \emph{strengthens} the absorption-quality association for sparse probing F1 from $r = -0.664$ to $r = -0.746$ ($p = 1.2 \times 10^{-9}$), a classical suppression effect.
Absorption fully mediates $L_0$'s effect on spurious correlation removal (SCR; direct effect $c' = -0.003$, n.s.; Sobel $z = 3.62$, $p = 2.9 \times 10^{-4}$), and the finding survives Rosenbaum sensitivity analysis with $\Gamma = 2.65$ for RAVEL TPP.
A first attempt to measure absorption beyond the first-letter spelling task---using five knowledge-hierarchy probes on GPT-2 Small---exposes a critical limitation of the standard dominance-based metric: shuffled-hierarchy controls produce 100\% absorption, indistinguishable from real hierarchies, while a cosine-calibrated variant detects 0\%.
Across 420 SAEs from SAEBench, a generalized additive model reveals a highly significant width--$L_0$ interaction ($p = 3.1 \times 10^{-15}$, $R^2 = 0.693$), with a transition zone at $L_0 \approx 7$--$14$ where absorption rises sharply; at $L_0 > 14$, absorption is low regardless of dictionary width.
These results establish absorption as an independent quality predictor, expose metric limitations for cross-domain measurement, and provide an actionable scaling map for practitioners selecting SAE hyperparameters.
\end{abstract}

%----------------------------------------------------------------------
\section{Introduction}
\label{sec:intro}
%----------------------------------------------------------------------

Sparse autoencoders (SAEs) decompose polysemantic neural network activations into overcomplete sparse bases of interpretable latents \citep{bricken2023monosemanticity,templeton2024scaling}.
The promise is mechanistic interpretability at scale: each SAE latent should correspond to a single concept, enabling researchers to audit and steer model behavior.
But SAEs suffer from a failure mode that directly undermines this promise.
\emph{Feature absorption} occurs when an SAE latent silently fails to fire because a more specific, co-occurring latent subsumes its role under the sparsity objective \citep{chanin2024absorption}.
On the first-letter spelling task---where a general ``starts with S'' latent is absorbed by a specific ``September'' latent---\citet{chanin2024absorption} measure 15--35\% absorption rates across hundreds of Gemma Scope SAEs, and every tested architecture (L1, TopK, JumpReLU, Gated) exhibits the phenomenon.
When a researcher steers model behavior by activating the ``starts with S'' latent, absorption means September-related inputs are silently excluded from the intervention.
Any downstream analysis that relies on SAE latent activations---circuit discovery, feature attribution, model editing---inherits these silent failures.

Three weaknesses block the field's ability to act on these measurements.

\paragraph{The confound problem.}
The strongest published evidence that absorption degrades SAE quality is a correlation of $r = -0.595$ across 54 Gemma Scope SAEs between absorption rate and downstream quality metrics \citep{chanin2024absorption,karvonen2025saebench}.
But this correlation was computed without controlling for $L_0$ (the expected number of active latents per input).
All high-absorption SAEs in that dataset are 1M-width with low $L_0$ (16--58); all low-absorption SAEs are 16k or 65k width with high $L_0$ (137--297).
If absorption is merely a proxy for $L_0$---which independently affects quality through feature hedging \citep{chanin2025hedging}---the entire absorption-quality narrative collapses into ``use SAEs with higher $L_0$.''
No prior work has applied formal confound control to this question.

\paragraph{The single-task problem.}
Every absorption measurement in the literature uses a single evaluation task: first-letter spelling.
At least five papers explicitly call out this limitation \citep{chanin2024absorption,karvonen2025saebench,korznikov2025ortsae,bussmann2025matryoshka,li2025atm}.
The first-letter task has an unusually clean, deterministic hierarchy---``September'' is always a member of ``starts with S''---that may systematically maximize the sparsity incentive for absorption.
Whether absorption occurs in fuzzier, real-world hierarchies (e.g., city-country, city-continent) is unknown.

\paragraph{The scaling problem.}
No study maps absorption jointly across SAE dictionary width and $L_0$.
Practitioners selecting SAE hyperparameters cannot determine whether there exists a region in (width, $L_0$) space where absorption is reliably low, or whether width and $L_0$ interact nonlinearly to produce absorption phase transitions.

This paper addresses all three weaknesses.

\textbf{Contribution 1: Confound resolution via epidemiological causal inference.}
We apply partial correlation analysis, Baron-Kenny mediation, and Rosenbaum sensitivity analysis to 48 Gemma Scope SAEs (the 54 in the prior study minus 6 lacking reported $L_0$)---methods standard in epidemiology and social science but never previously applied to SAE evaluation.
The result reverses the concern that absorption might be an $L_0$ epiphenomenon: after controlling for $\log(L_0)$, the partial correlation between absorption and sparse probing F1 strengthens from $r = -0.664$ to $r = -0.746$, a classical suppression effect where $L_0$ was partially masking absorption's true impact.
Mediation analysis establishes absorption as the primary pathway through which $L_0$ affects SCR (full mediation) and TPP (proportion mediated $= 0.54$).
Rosenbaum sensitivity bounds show that the TPP result can withstand a hidden confounder with a 2.65:1 odds ratio (Figure~\ref{fig:partial_corr}; Figure~\ref{fig:mediation}).

\textbf{Contribution 2: Cross-domain absorption measurement on knowledge hierarchies.}
We provide the first measurement of absorption beyond first-letter spelling, using five probe types over 3,552 cities from the RAVEL dataset on GPT-2 Small (124M parameters; used as a fallback due to Gemma 2B access restrictions).
The attempt exposes a critical methodological limitation: the standard dominance-based absorption metric \citep{chanin2024absorption} produces absorption rates of 11.3--96.2\% across knowledge domains and layers, but shuffled-hierarchy controls---where city-attribute mappings are randomized---show 100\% absorption.
A cosine-calibrated variant that requires alignment between the dominant feature and the probe direction detects 0\% absorption across all thresholds (Figure~\ref{fig:crossdomain}).

\textbf{Contribution 3: Absorption scaling surface with significant interaction structure.}
We construct the first empirical absorption phase surface across 420 SAEs from SAEBench (Gemma 2 2B), spanning dictionary widths from 2,304 to 1,048,576 and $L_0$ from 9.3 to 8,277.
A generalized additive model (GAM) with a tensor interaction term yields $R^2 = 0.693$, substantially outperforming both the additive model ($R^2 = 0.620$) and the linear baseline ($R^2 = 0.488$).
The interaction term is highly significant ($p = 3.1 \times 10^{-15}$): absorption cannot be predicted from width or $L_0$ independently.
Gradient analysis identifies a transition zone at $L_0 \approx 7$--$14$, where absorption rises sharply.
At $L_0 > 14$, absorption is low regardless of width; at $L_0 < 7$, absorption increases dramatically as width scales from 16k to 1M (Figure~\ref{fig:scaling}).

\textbf{Contribution 4: Taxonomy correction reveals inflated baseline.}
Frequency-matched correction of the original \citet{chanin2024absorption} absorption taxonomy on GPT-2 Small reveals that the reported 92.3\% comprehensive absorption rate was an artifact of feature specificity: the corrected rate drops to 19.2\% when proper non-letter-context comparison baselines are used, while the Chanin false-negative-based metric independently validates absorption in 73.1\% of letters.

Together, these contributions transform feature absorption from a narrowly validated observation on a single task into a rigorously characterized phenomenon with quantified causal status, documented metric limitations, and an actionable scaling map.

Section~\ref{sec:background} contextualizes the confound problem and prior mitigation work.
Sections~\ref{sec:method} and~\ref{sec:experiments} present methods and results for each contribution in sequence, followed by interpretation (Section~\ref{sec:discussion}) and practical recommendations (Section~\ref{sec:conclusion}).

%----------------------------------------------------------------------
\section{Background and Related Work}
\label{sec:background}
%----------------------------------------------------------------------

\subsection{Sparse Autoencoders for Mechanistic Interpretability}

Given an input activation $\mathbf{x} \in \mathbb{R}^d$, an SAE computes latent activations $\mathbf{z} \in \mathbb{R}^m$ (where $m \gg d$) and reconstructs $\hat{\mathbf{x}} = W_d \mathbf{z} + \mathbf{b}_d$.
The sparsity objective ensures that $L_0$ (the expected number of active latents per input) is small relative to $m$, so that each active latent can be interpreted as a single semantic direction.

Five SAE architectures have emerged since 2023, each addressing a different training limitation.
\citet{bricken2023monosemanticity} demonstrated the first large-scale monosemantic decomposition on a 1-layer transformer using L1-penalized SAEs.
\citet{gao2024topk} proposed TopK SAEs that select exactly $k$ highest-activating latents per input, yielding clean scaling laws on GPT-4 with up to 16M latents.
\citet{rajamanoharan2024gated} introduced Gated SAEs, which decouple feature detection from magnitude estimation to reduce L1-induced shrinkage bias.
\citet{rajamanoharan2024jumprelu} proposed JumpReLU SAEs, which train $L_0$ directly via the straight-through estimator and achieve state-of-the-art reconstruction fidelity on Gemma 2 9B.
Google DeepMind released Gemma Scope \citep{lieberum2024gemmascope}, comprising 400+ open-source JumpReLU SAEs across all layers of Gemma 2 2B, 9B, and 27B with dictionary widths from 1k to 1M---the primary evaluation target for absorption research.

Evaluation has standardized around SAEBench \citep{karvonen2025saebench}, which scores 200+ open-source SAEs on eight metrics: sparse probing, RAVEL disentanglement (TPP), spurious correlation removal (SCR), unlearning, and absorption, among others.
A key SAEBench finding motivates our work: proxy metrics such as CE loss recovered and $L_0$ do not reliably predict practical downstream performance \citep[Table~3]{karvonen2025saebench}, raising the question of whether absorption---a structural failure mode rather than a proxy metric---is a better quality predictor.

\subsection{Feature Absorption: Definition and Prior Measurements}

\citet{chanin2024absorption} formalized feature absorption.
The phenomenon occurs when a parent SAE latent $j$ (e.g., ``starts with S'') fails to fire on tokens where a more specific child latent $c$ (e.g., ``September'') activates instead, because the SAE's sparsity objective favors the child's single-latent encoding over the parent-plus-child two-latent encoding.
Their measurement protocol proceeds in four steps:
\begin{enumerate}
\item Train linear probes to identify ground-truth feature directions.
\item Use $k$-sparse probing to find feature splits.
\item Identify false-negative tokens (where all split latents for a feature fail to fire but the probe correctly classifies the input).
\item Apply integrated-gradients ablation to attribute false negatives to specific absorbing latents. Absorption is detected when the highest-ablation-effect latent has cosine similarity exceeding $\tau_{\text{cos}} = 0.025$ with the parent probe direction and dominance ratio exceeding $\tau_{\text{dom}} = 1.0$ over the second-highest.
\end{enumerate}

On the first-letter spelling task across mid-layer Gemma Scope SAEs, \citet{chanin2024absorption} report absorption rates of 15--35\%, with the rate increasing at wider dictionaries and lower $L_0$.
Their toy model proves that absorption is loss-optimal when the sparsity penalty exceeds $\sin^2(\theta_{p,c})$, where $\theta_{p,c}$ is the decoder angle between parent and child directions.
Absorption appears in every tested SAE architecture (L1, TopK, JumpReLU) and across multiple model families (Gemma 2, Llama 3.2, Qwen2).

\citet{tian2025sensitivity} frame absorption as a special case of poor feature \emph{sensitivity}: a latent that activates selectively on its target concept but fails on similar inputs has low sensitivity.
Their scalable evaluation, which covers thousands of features across multiple SAE families, reveals that many features rated as interpretable by activation-example inspection have poor recall, consistent with the partial absorption phenomenon.

Two properties of the existing measurement base limit the field's understanding.
All absorption measurements use a single evaluation task---the first-letter spelling task---a deterministic hierarchy with an unusually sharp parent-child structure.
Whether absorption rates generalize to fuzzier, semantically richer hierarchies (knowledge taxonomies, safety-relevant features) is unknown, a limitation explicitly noted by \citet{chanin2024absorption} and at least four subsequent papers \citep{karvonen2025saebench,korznikov2025ortsae,bussmann2025matryoshka,li2025atm}.
No study has tested whether absorption scores predict downstream SAE quality after controlling for confounds, leaving the assumed causal chain (less absorption implies better interpretability) empirically unvalidated.

\subsection{Mitigation Approaches}

Multiple architectural interventions reduce absorption, though no unified account explains why they work.

\citet{bussmann2025matryoshka} proposed Matryoshka SAEs, which train nested dictionaries of increasing width simultaneously.
The nested structure allocates general features to smaller inner dictionaries and specific features to larger outer ones, directly reducing the parent-child competition that drives absorption.
Matryoshka SAEs achieve the best absorption scores in SAEBench while maintaining competitive reconstruction, though inner dictionary levels suffer from feature hedging \citep{chanin2025hedging}.

\citet{korznikov2025ortsae} enforce pairwise orthogonality on decoder columns via a cosine similarity penalty (OrtSAE), reducing absorption by 65\% relative to standard SAEs with linear computational overhead.
Orthogonality decreases $\cos(\mathbf{p}, \mathbf{d}_j)$ between parent probe directions and child decoder directions, directly targeting the geometric condition in the \citet{chanin2024absorption} toy model.

\citet{li2025atm} introduced Adaptive Temporal Masking (ATM), which dynamically scores per-latent importance based on activation magnitude, frequency, and reconstruction contribution.
ATM achieves the lowest reported absorption score: 0.0068 on Gemma 2 2B, compared to 0.1402 for TopK and 0.0114 for JumpReLU.

\citet{narayanaswamy2026masked} proposed masked regularization, which randomly masks high-frequency tokens during SAE training to disrupt the co-occurrence patterns that enable absorption.
The approach improves out-of-distribution robustness across SAE architectures, though quantitative absorption reduction numbers are not reported using the standard Chanin metric.

A unified theoretical framework by \citet{wright2025framework} casts all sparse dictionary learning methods as a piecewise biconvex optimization problem and identifies stable partial minima where absorbed features are trapped.
The proposed remedy, feature anchoring, restores identifiability in synthetic benchmarks but has not been validated for absorption reduction on real LLM SAEs.

These diverse mechanisms---nesting, orthogonality, temporal masking, token masking, anchoring---all reduce absorption but operate through different pathways.
Our scaling surface analysis (Section~\ref{sec:scaling}) provides complementary evidence by mapping which regions of the (width, $L_0$) hyperparameter space exhibit high absorption across the 420 SAEs in the SAEBench collection.

\subsection{The Unresolved Confound Problem}

The most consequential open question about absorption is whether it is a genuine causal driver of SAE quality degradation or an epiphenomenon of correlated hyperparameters.

\citet{chanin2024absorption} report that absorption correlates with downstream quality ($r = -0.595$ across 54 Gemma Scope SAEs), but their analysis does not control for $L_0$.
In the Gemma Scope collection, all high-absorption SAEs are 1M width with low $L_0$ (16--58), while all low-absorption SAEs are 16k or 65k width with high $L_0$ (137--297).
Prior partial correlations controlled for $\log(\text{width})$ and layer but not $L_0$, leaving open the possibility that absorption is simply a proxy for sparsity level---in which case the entire mitigation research program is misdirected.

The confound is sharpened by \citet{chanin2025hedging}, who show that incorrect $L_0$ (which is common in open-source SAEs) causes feature hedging, a distinct failure mode where correlated features are merged into a single latent.
Feature hedging and absorption both manifest as features failing to fire where they should, but they have different causes: hedging arises from insufficient dictionary capacity, while absorption arises from hierarchical feature structure.
No study has systematically disentangled the two in observational SAE data.

No prior work applies formal causal inference methods to SAE evaluation.
The SAE community relies on bivariate or partially controlled correlations between architectures, without the mediation analysis, propensity matching, or sensitivity analysis that are standard in epidemiological and social science research for establishing causal claims from observational data.
Our Phase~1 analysis (Section~\ref{sec:confound}) introduces these methods to SAE evaluation.
Our Phase~2 (Section~\ref{sec:crossdomain}) tests whether the standard absorption metric transfers to knowledge hierarchies.
Our Phase~3 (Section~\ref{sec:scaling}) constructs the first joint (width, $L_0$) absorption scaling surface with formal interaction testing across 420 SAEs.

%----------------------------------------------------------------------
\section{Method}
\label{sec:method}
%----------------------------------------------------------------------

This study comprises four analysis phases: (1) confound resolution via epidemiological causal inference applied to 48 Gemma Scope SAEs, (2) cross-domain absorption measurement on knowledge hierarchies using GPT-2 Small, (3) absorption scaling surface construction across 420 SAEs with formal interaction testing, and (4) taxonomy correction for the original absorption classification.
All analyses are training-free, operating on existing pre-trained SAEs and precomputed SAEBench results.
Figure~\ref{fig:mediation} illustrates the mediation path model central to Phase~1.

\subsection{Confound Resolution via Epidemiological Causal Inference}
\label{sec:confound}

\subsubsection{Dataset}

The confound resolution analysis uses 48 Gemma Scope SAEs (Gemma 2 2B) from SAEBench \citep{karvonen2025saebench}.
Six SAEs from the original 54 that lack reported $L_0$ values are excluded.
Each SAE has a precomputed absorption score \citep{chanin2024absorption} and four downstream quality metrics: sparse probing F1 ($\text{SP-F1}$), spurious correlation removal ($\text{SCR}$), RAVEL true positive proportion ($\text{TPP}$), and unlearning ($\text{UL}$).
SAE dictionary widths span 16,384 to 1,048,576 and $L_0$ ranges from 9 to 297.
All 48 SAEs share a single architecture class (JumpReLU), so architecture is dropped as a covariate.

\subsubsection{Step 1: $L_0$ as Covariate (Go/No-Go)}

We compute partial Pearson correlations between absorption rate $R_{\text{abs}}$ and each quality metric, controlling for $\log(\text{width})$, layer, and $\log(L_0)$:
\[
r_{\text{partial}}(R_{\text{abs}}, Q_i \mid \log m, \ell, \log L_0)
\]
where $Q_i \in \{\text{SP-F1}, \text{SCR}, \text{TPP}, \text{UL}\}$, $m$ is dictionary width, and $\ell$ is layer index.
This extends the prior analysis by \citet{chanin2024absorption} (which controlled for $\log m$ and $\ell$ but not $L_0$) by adding the critical confound identified in Section~\ref{sec:background}.

\textbf{Go/no-go criterion}: At least one quality metric must retain $|r_{\text{partial}}| > 0.2$ after $L_0$ control.
If all four drop below this threshold, the absorption-quality association is an artifact of the width/$L_0$ confound, and H1 is falsified.

Collinearity is assessed via variance inflation factors (VIF).
With $\text{VIF}(\log L_0) = 1.09$ and $\text{VIF}(\log m) = 1.08$, multicollinearity is not a concern.

\subsubsection{Step 2: Width-Stratified Analysis}

SAEs are partitioned into three width strata: 16k ($n = 15$), 65k ($n = 15$), and 1M ($n = 18$).
Within each stratum, we compute Spearman rank correlations $\rho$ between $R_{\text{abs}}$ and each quality metric, with BCa bootstrap 95\% confidence intervals (10,000 resamples).
This tests whether the absorption-quality association holds within fixed width or is driven entirely by cross-width variation.

\subsubsection{Step 3: Baron-Kenny Mediation Analysis}

We test the causal path $\log(L_0) \to R_{\text{abs}} \to Q_i$ using the Baron-Kenny four-step procedure, controlling for $\log m$ and $\ell$:
\begin{enumerate}
\item \textbf{Equation 1 (path $a$)}: $R_{\text{abs}} = \beta_0 + a \cdot \log(L_0) + \gamma_1 \log m + \gamma_2 \ell + \epsilon$
\item \textbf{Equation 2 (path $b$ and direct effect $c'$)}: $Q_i = \beta_0 + b \cdot R_{\text{abs}} + c' \cdot \log(L_0) + \gamma_1 \log m + \gamma_2 \ell + \epsilon$
\item \textbf{Equation 3 (total effect $c$)}: $Q_i = \beta_0 + c \cdot \log(L_0) + \gamma_1 \log m + \gamma_2 \ell + \epsilon$
\item \textbf{Indirect effect}: $ab = c - c'$
\end{enumerate}

Statistical significance of the indirect effect is assessed via the Sobel test and 10,000-resample bootstrap confidence intervals.
Full mediation is declared when all four Baron-Kenny conditions are met and $c'$ is non-significant; partial mediation when $c'$ is reduced but remains significant.

\subsubsection{Step 4: Rosenbaum Sensitivity Analysis}

High-absorption ($R_{\text{abs}} > 0.3$) and low-absorption ($R_{\text{abs}} < 0.1$) SAEs are matched on width (exact) and $L_0$ (nearest-neighbor Mahalanobis distance matching).
For each matched pair, quality differences are tested via the Wilcoxon signed-rank test.
The Rosenbaum sensitivity parameter $\Gamma$ is computed---the odds ratio of a hypothetical hidden confounder at which the matched-pair result becomes non-significant.

Two matching strategies are employed: propensity score matching (6 pairs) and Mahalanobis matching (17 pairs).
Within-width matching strategies (median split: 23 pairs; tertile: 16 pairs) serve as additional controls.

\subsubsection{Step 5: SCR Suppression Variable Diagnosis}

To identify which covariate produces the suppression effect on SCR, covariates are added sequentially: width-only, layer-only, architecture-only, $L_0$-only.
The partial correlation between absorption and SCR is tracked at each step, isolating which variable generates the change from the bivariate $r = -0.431$ to the full partial $r$.

\subsection{Cross-Domain Absorption Measurement}
\label{sec:crossdomain}

\subsubsection{Model and SAE}

Cross-domain experiments use GPT-2 Small (124M parameters; \citealp{radford2019language}) with the \texttt{gpt2-small-res-jb} SAE (24,576 latents), accessed via SAELens \citep{bloom2024saelens} and TransformerLens \citep{nanda2022transformerlens}.
Gemma 2 2B was the intended target, but HuggingFace access restrictions at the time of experimentation necessitated the GPT-2 Small fallback.
This model change means that H2 is tested on a different model-SAE pair than originally planned, with a substantially smaller model (124M vs.\ 2B) and smaller SAE dictionary (24k vs.\ 16k--1M).
We address the implications of this change in Section~\ref{sec:limitations}.

\subsubsection{Probe Training}

Logistic regression probes are trained on residual stream activations at \texttt{hook\_resid\_pre} (matching the SAE input hook point) for five attribute types across three layers (5, 8, 11---spanning GPT-2 Small's early, mid, and late layers):

\begin{table}[h]
\centering
\small
\begin{tabular}{@{}llll@{}}
\toprule
Probe Type & Granularity & Classes & Training Samples \\
\midrule
Country (US) & Binary & US / non-US & 3,552 cities \\
Language (English) & Binary & English / non-English & 3,552 cities \\
Continent & 6-class & 6 continents & 3,552 cities \\
Country top-10 & 10-class & 10 most frequent countries & 1,523 cities \\
Language top-10 & 10-class & 10 most frequent languages & 2,297 cities \\
\bottomrule
\end{tabular}
\end{table}

Probes are trained with scikit-learn \texttt{LogisticRegression(max\_iter=1000, C=1.0)} using 80/20 stratified train/test splits, averaged over 3 random seeds (42, 123, 456).
Prompt template: \texttt{"The city of \{City\} is located in"} with activations extracted at the last token position.
A quality gate of 85\% accuracy is applied to binary probes; multiclass probes are included with explicit accuracy reporting.

\subsubsection{Absorption Measurement (Adapted Chanin Metric)}

The absorption measurement adapts the \citet{chanin2024absorption} protocol to knowledge hierarchies via four steps:
\begin{enumerate}
\item \textbf{$k$-sparse probing}: For each attribute value, identify $k$ split latents---SAE latents with selectivity $\geq 3.0$ (activation rate $\geq$ 5\% on the target class and $\geq 3\times$ higher than on other classes).
\item \textbf{False-negative identification}: Find tokens in $\text{FN}(f)$---inputs where all $k$ split latents for feature $f$ fail to activate but the linear probe correctly classifies the token.
\item \textbf{Dominance-based detection}: On false-negative tokens, the SAE latent with the highest activation magnitude is identified. A token is classified as absorbed if this dominant latent's activation exceeds the second-highest by a dominance ratio $\tau_{\text{dom}} \geq 1.0$.
\item \textbf{Cosine-calibrated detection}: As an alternative, absorption is declared only when the dominant latent's decoder direction $\mathbf{d}_j$ has $\cos(\mathbf{p}, \mathbf{d}_j) > \tau_{\text{cos}}$ with the probe direction $\mathbf{p}$.
\end{enumerate}

\subsubsection{Controls}

Three control conditions validate the absorption metric:
\begin{itemize}
\item \textbf{Shuffled hierarchy}: City-attribute mappings are randomized (5 trials per probe per layer), destroying real hierarchical structure.
\item \textbf{Random probe direction}: Random unit vectors replace trained probe directions (5 trials per layer).
\item \textbf{First-letter baseline}: The standard \citet{chanin2024absorption} first-letter spelling absorption measurement on the same SAE.
\end{itemize}

\subsubsection{Threshold Sweep}

Absorption rates are computed across a grid of detection thresholds: cosine similarity $\tau_{\text{cos}} \in \{0.05, 0.10, 0.15, 0.20, 0.30\}$ and dominance ratio $\tau_{\text{dom}} \in \{1.0, 2.0\}$.

\subsection{Absorption Scaling Surface}
\label{sec:scaling}

\subsubsection{Dataset}

The scaling surface uses 420 SAEs from SAEBench spanning 9 release families of Gemma 2 2B SAEs.
Dictionary widths range from 2,304 to 1,048,576 ($\log_2$ range: 11.2--20.0); $L_0$ ranges from 9.3 to 8,277 ($\log_2$ range: 3.2--13.0).
Architecture breakdown: 360 standard (L1), 54 JumpReLU, 6 unclassified.
Three layers are represented (5, 12, 19), with 140 SAEs per layer.

\subsubsection{GAM Fit and Model Comparison}

Three regression models of increasing flexibility are compared:

\textbf{Linear model}:
\[
R_{\text{abs}} = \beta_0 + \beta_1 \log_2 m + \beta_2 \log_2 L_0 + \beta_3 \ell + \epsilon
\]

\textbf{Additive GAM}:
\[
R_{\text{abs}} = s(\log_2 m) + s(\log_2 L_0) + \beta_3 \ell + \epsilon
\]

\textbf{Interaction GAM}:
\[
R_{\text{abs}} = s(\log_2 m) + s(\log_2 L_0) + \text{ti}(\log_2 m, \log_2 L_0) + \beta_3 \ell + \epsilon
\]
where $s(\cdot)$ denotes smooth univariate spline terms and $\text{ti}(\cdot, \cdot)$ is a tensor interaction term that captures joint dependence of absorption on width and $L_0$ beyond their additive contributions.

\subsubsection{Phase Boundary Detection}

The gradient magnitude of the fitted GAM surface is computed on a dense grid in $(\log_2 m, \log_2 L_0)$ space.
Ridges---regions of maximal gradient magnitude---are identified as candidate transition zones where absorption changes rapidly.
Ridge points are defined as grid cells where gradient magnitude exceeds the 70th percentile of the overall distribution.

\subsection{Taxonomy Correction}
\label{sec:taxonomy}

\subsubsection{Motivation}

The original \citet{chanin2024absorption} absorption taxonomy classifies 92.3\% of letters as exhibiting some form of absorption on GPT-2 Small.
The Type~II classification (88.5\% of letters) relies on a magnitude ratio comparing a parent feature's activation in letter-specific contexts versus comparison contexts.
For letters with zero comparison tokens, the original implementation falls back to a global mean-when-active baseline, potentially inflating the Type~II rate.

\subsubsection{Frequency-Matched Correction}

For each of the 26 letters, we apply a three-strategy comparison baseline:
\begin{itemize}
\item \textbf{Strategy A (preferred)}: Non-letter-context activations of the same parent feature on WikiText-103 (7,991 sentences).
\item \textbf{Strategy B (fallback)}: Global when-active baseline from the WikiText-103 corpus when Strategy A yields insufficient data.
\item \textbf{Strategy C (independent validation)}: The Chanin absorption rate (fraction of false-negative tokens per letter).
\end{itemize}

The Type~II classification threshold remains at magnitude ratio $< 0.5$.
The corrected combined rate is computed with BCa bootstrap 95\% confidence intervals (10,000 resamples).

\subsection{Pre-Registered Hypotheses and Falsification Criteria}

\textbf{H1} (Absorption-Quality Causal Chain): After controlling for $\log(L_0)$, the partial correlation between absorption and at least one quality metric retains $|r_{\text{partial}}| > 0.2$ ($p < 0.05$).
Absorption mediates ${>}\,30\%$ of $L_0$'s total effect on at least one downstream quality metric (bootstrap CI excluding 0).

\emph{Falsified if}: All four partial correlations drop below $|0.2|$, all width-stratified correlations are non-significant, mediation indirect effect bootstrap CIs include 0, and $\Gamma < 1.2$.

\textbf{H2} (Cross-Domain Absorption Existence): Absorption rate on knowledge hierarchies exceeds 10\% and exceeds $3\times$ the shuffled-hierarchy baseline for at least one domain-layer combination.

\emph{Falsified if}: No hierarchy shows absorption exceeding $3\times$ the shuffled baseline after the probe quality gate.

\textbf{H3} (Absorption Scaling Surface): The tensor interaction term $\text{ti}(\log_2 m, \log_2 L_0)$ in the GAM is significant ($p < 0.05$).

\emph{Falsified if}: The interaction term is non-significant ($p > 0.10$) and the additive GAM explains absorption as well as the interaction GAM.

\textbf{H4} (Taxonomy Correction): The original 92.3\% comprehensive absorption rate is inflated by the global baseline fallback.
Frequency-matched correction will reduce the rate.

\emph{Falsified if}: The corrected rate is within 5 percentage points of the original 92.3\%.

\subsection{Software and Reproducibility}

All analyses use Python 3.12 with numpy, scipy, pandas, statsmodels, pingouin (partial correlations), scikit-learn (probes and matching), pyGAM (generalized additive models), and matplotlib/seaborn (visualization).
Model access uses TransformerLens $\geq$ 2.0 and SAELens $\geq$ 4.0.
Random seeds are fixed at 42 (primary), 123, and 456 (replication seeds).
All threshold sweeps use pre-registered values.

%----------------------------------------------------------------------
\section{Experiments}
\label{sec:experiments}
%----------------------------------------------------------------------

\subsection{Confound Resolution: Absorption Survives $L_0$ Control (H1)}

\subsubsection{Setup}

We analyze 48 Gemma Scope SAEs (Gemma 2 2B) with known $L_0$ values from the SAEBench dataset.
All 48 remaining SAEs share the same architecture class (JumpReLU), so architecture class is dropped as a covariate.
Covariates are $\log(\text{width})$, layer, and $\log(L_0)$.
Four quality metrics are evaluated: Sparse Probing F1 (SP-F1, $n = 48$), Spurious Correlation Removal (SCR, $n = 43$), RAVEL True Positive Proportion (TPP, $n = 48$), and Unlearning (UL, $n = 34$).

\subsubsection{Go/No-Go: $L_0$ as Covariate}

The critical test is whether the absorption-quality partial correlation survives the addition of $\log(L_0)$.
Table~\ref{tab:partial_corr} reports the results.

\begin{table}[t]
\caption{Partial correlations before and after $L_0$ control. Three of four quality metrics retain $|r| > 0.2$ after controlling for $L_0$. The sparse probing suppression effect (bold) is the central finding.}
\label{tab:partial_corr}
\centering
\small
\begin{tabular}{@{}lccccc@{}}
\toprule
Quality Metric & $n$ & Bivariate $r$ & Partial $r$ (no $L_0$) & Partial $r$ (with $L_0$) & $p$ (with $L_0$) \\
\midrule
\textbf{Sparse Probing F1} & 48 & $-0.587$ & $-0.664$ & $\mathbf{-0.746}$ & $1.2 \times 10^{-9}$ \\
SCR & 43 & $-0.449$ & $-0.692$ & $-0.570$ & $6.6 \times 10^{-5}$ \\
RAVEL TPP & 48 & $-0.471$ & $-0.488$ & $-0.331$ & $0.022$ \\
Unlearning & 34 & $-0.182$ & $-0.082$ & $-0.123$ & $0.487$ \\
\bottomrule
\end{tabular}
\end{table}

Three of four quality metrics retain $|r_{\text{partial}}| > 0.2$ after controlling for $L_0$ (SP-F1, SCR, TPP).
The go/no-go criterion is met.

The sparse probing result contains a suppression effect: the partial correlation \emph{strengthened} from $r = -0.664$ to $r = -0.746$ when $L_0$ is added as a covariate.
$L_0$ shares variance with absorption that does not predict sparse probing quality.
Controlling for $L_0$ removes this non-predictive variance and unmasks the true absorption-quality relationship.
Figure~\ref{fig:partial_corr} illustrates this pattern across all four metrics.

\begin{figure}[t]
\centering
\includegraphics[width=0.85\linewidth]{figures/fig1_partial_correlations.pdf}
\caption{Grouped bar chart comparing partial correlations between absorption and four downstream quality metrics, with and without $L_0$ as a covariate. The dashed line marks the $|r| = 0.2$ go/no-go threshold. The sparse probing suppression effect---where the partial correlation strengthens from $-0.664$ to $-0.746$ upon $L_0$ control---is the central finding.}
\label{fig:partial_corr}
\end{figure}

Collinearity diagnostics confirm that the covariates are sufficiently independent: VIF values are 1.08 ($\log$ width), 1.01 (layer), and 1.09 ($\log L_0$), all below the standard threshold of 5.

For SCR, the partial correlation weakens from $r = -0.692$ (without $L_0$) to $r = -0.570$ (with $L_0$), a reduction of 17.6\%.
The SCR suppression decomposition analysis reveals that layer is the primary suppressor variable: controlling for layer alone shifts the bivariate SCR correlation from $r = -0.449$ to $r = -0.836$, the largest single-covariate effect.

For unlearning, the association is non-significant in all analyses ($p = 0.487$), consistent with the expectation that unlearning captures different aspects of SAE quality.

\subsubsection{Width-Stratified Analysis}

Within-width correlations test whether the absorption-quality link persists after removing cross-width variation.
Table~\ref{tab:stratified} reports Spearman $\rho$ within each width stratum.

\begin{table}[t]
\caption{Width-stratified Spearman correlations (BCa bootstrap 95\% CI). The 1M stratum shows the strongest trends. No individual stratum achieves a CI excluding zero, reflecting limited per-stratum sample sizes.}
\label{tab:stratified}
\centering
\small
\begin{tabular}{@{}ccll l@{}}
\toprule
Width & $n$ & SP-F1 $\rho$ [CI] & SCR $\rho$ [CI] & TPP $\rho$ [CI] \\
\midrule
16k & 15 & $-0.236$ [$-0.722, 0.336$] & $0.025$ [$-0.556, 0.575$] & $-0.296$ [$-0.786, 0.311$] \\
65k & 15 & $0.264$ [$-0.499, 0.763$] & $0.021$ [$-0.596, 0.520$] & $-0.100$ [$-0.672, 0.455$] \\
1M & 18 & $-0.480$ [$-0.829, 0.076$] & $-0.549$ [$-0.872, 0.179$] & $-0.399$ [$-0.732, 0.118$] \\
\bottomrule
\end{tabular}
\end{table}

No individual stratum achieves a 95\% CI excluding zero, reflecting the limited per-stratum sample sizes ($n = 15$--$18$).
The 1M stratum (widest $L_0$ range: 9--207) shows the strongest trends: $\rho = -0.480$ ($p = 0.044$) for sparse probing and $\rho = -0.549$ ($p = 0.052$) for SCR.
The pooled Spearman analysis across all strata confirms 3/4 metrics with bootstrap CI excluding zero: SP-F1 ($\rho = -0.455$, $p = 0.001$), SCR ($\rho = -0.376$, $p = 0.013$), TPP ($\rho = -0.524$, $p = 1.3 \times 10^{-4}$).

\subsubsection{Baron-Kenny Mediation Analysis}

We test the causal path $\log(L_0) \to \text{Absorption} \to \text{Quality}$, controlling for $\log(\text{width})$ and layer, with 10,000 bootstrap resamples.

\begin{table}[t]
\caption{Mediation analysis results. SCR shows full Baron-Kenny mediation (bold): the direct effect $c'$ is non-significant, and the indirect effect bootstrap CI excludes zero. All path coefficients are standardized.}
\label{tab:mediation}
\centering
\footnotesize
\setlength{\tabcolsep}{3.5pt}
\begin{tabular}{@{}lcccccc@{}}
\toprule
Metric & Type & $a$ (std) & $b$ (std) & $c'$ ($p$) & $ab$ [95\% CI] & Sobel $z$ ($p$) \\
\midrule
SP-F1 & Indir.\ only$^\dagger$ & $-0.47$ & $-0.73$ & $-0.27$ (.001) & .015 [.007, .028] & 4.08 ($4.4{\times}10^{-5}$) \\
\textbf{SCR} & \textbf{Full} & $\mathbf{-0.54}$ & $\mathbf{-0.46}$ & $\mathbf{-0.03}$ \textbf{(.71)} & \textbf{.025 [.007, .048]} & \textbf{3.62} ($\mathbf{2.9{\times}10^{-4}}$) \\
TPP & Full & $-0.47$ & $-0.31$ & $0.13$ (.25) & .003 [$10^{-5}$, .007] & 2.08 (.037) \\
UL & None & $-0.68$ & $-0.09$ & $-0.07$ (.62) & .007 [$-.028$, .035] & 0.66 (.51) \\
\bottomrule
\end{tabular}
\vspace{0.5em}

{\footnotesize $^\dagger$Per \citet{zhao2010reconsidering}: the indirect effect is significant but the total effect $c$ is non-significant ($p = 0.45$), precluding classical Baron-Kenny mediation. The significant indirect effect with non-significant total effect indicates \emph{indirect-only} mediation (also called suppression-mediation): the direct effect $c' = -0.270$ ($p = 0.001$) is negative and significant, while the indirect effect $ab = 0.015$ is positive, reflecting opposite-sign pathways. $L_0$ directly degrades quality through a non-absorption mechanism, but simultaneously reduces absorption, which improves quality. These opposing paths cancel in the total effect.}
\end{table}

SCR shows full Baron-Kenny mediation: all four conditions are met, the direct effect $c' = -0.029$ is non-significant ($p = 0.71$), and the indirect effect bootstrap CI excludes zero.
Absorption fully mediates the effect of $L_0$ on SCR.
Figure~\ref{fig:mediation} illustrates the path diagram with standardized coefficients.

\begin{figure}[t]
\centering
\includegraphics[width=0.75\linewidth]{figures/fig2_mediation_path.pdf}
\caption{Baron-Kenny mediation path diagram for SCR. Standardized coefficients: path $a$ ($\log(L_0)$ to absorption) $= -0.543$; path $b$ (absorption to SCR) $= -0.457$; direct effect $c' = -0.029$ (n.s.). Indirect effect $= 0.025$, 95\% CI $[0.007, 0.048]$. Absorption fully mediates $L_0$'s effect on SCR quality.}
\label{fig:mediation}
\end{figure}

TPP also meets full mediation criteria by the Baron-Kenny procedure (direct effect $p = 0.25$, indirect CI excludes zero), with proportion mediated $= 0.54$.

\subsubsection{Rosenbaum Sensitivity Analysis}

Five matching strategies are compared.
Mahalanobis matching (17 pairs, matching on $\log$ width, layer, and $\log L_0$) achieves the strongest results.

\begin{table}[t]
\caption{Rosenbaum sensitivity analysis by matching strategy. The $\Gamma$ parameter indicates the odds ratio of a hidden confounder needed to nullify the result. Higher $\Gamma$ indicates greater robustness.}
\label{tab:rosenbaum}
\centering
\small
\begin{tabular}{@{}lccccc@{}}
\toprule
Strategy & $n$ Pairs & SP-F1 $\Gamma$ & SCR $\Gamma$ & TPP $\Gamma$ \\
\midrule
Exact width, NN $L_0$ & 4 & 1.0 & 1.0 & 1.0 \\
Within-width median & 23 & 1.0 & 1.0 & 1.0 \\
Propensity score & 6 & 1.8 & 1.0 & 1.8 \\
\textbf{Mahalanobis} & \textbf{17} & \textbf{1.85} & \textbf{1.15} & \textbf{2.65} \\
Tertile within-width & 16 & 1.0 & 1.0 & 1.0 \\
\bottomrule
\end{tabular}
\end{table}

Under Mahalanobis matching, the TPP result achieves $\Gamma = 2.65$: the finding would remain significant even if a hidden confounder altered the odds of treatment assignment by a factor of 2.65:1.
Sparse probing reaches $\Gamma = 1.85$ (moderate robustness).
Exact-width and within-width matching strategies (4--23 pairs) fail to detect significant differences, consistent with the width-stratified null: the absorption-quality effect is strongest in cross-width comparisons.

\subsection{Cross-Domain Absorption: Metric Limitation Exposed (H2)}
\label{sec:crossdomain_results}

\subsubsection{Setup}

We use GPT-2 Small (124M parameters) with the \texttt{gpt2-small-res-jb} SAE (24,576 latents) at layers 5, 8, and 11.
The evaluation dataset contains 3,552 cities from the RAVEL knowledge dataset.
Binary probes achieve 84--94\% accuracy; multiclass probes range from 63.3\% to 81.0\%.

\subsubsection{Dominance-Based Absorption Rates}

Dominance-based detection (selectivity threshold $= 3$, dominance threshold $= 1.0$) produces absorption rates of 11.3--96.2\% across all domain-layer combinations.
Figure~\ref{fig:crossdomain} summarizes the cross-domain absorption rates with the shuffled control.

\begin{figure}[t]
\centering
\includegraphics[width=0.85\linewidth]{figures/fig3_crossdomain_absorption.pdf}
\caption{Dominance-based absorption rates across five knowledge-hierarchy probe types, averaged over layers. The dashed red line marks the shuffled-hierarchy control rate of 100\%. The metric does not discriminate real from randomized hierarchies.}
\label{fig:crossdomain}
\end{figure}

\begin{table}[t]
\caption{Cross-domain absorption rates by domain and layer. The $R_{\text{abs}}$ column equals the FN Rate in every row because at $\tau_{\text{dom}} = 1.0$, any non-split feature with nonzero activation satisfies the dominance criterion.}
\label{tab:crossdomain}
\centering
\small
\begin{tabular}{@{}lccccc@{}}
\toprule
Probe Type & Layer & Accuracy & $n$ Split & FN Rate & $R_{\text{abs}}$ \\
\midrule
Country binary US & 5 & 0.929 & 6 & 0.886 & 0.886 \\
Country binary US & 8 & 0.933 & 18 & 0.537 & 0.537 \\
Country binary US & 11 & 0.939 & 51 & 0.113 & \textbf{0.113} \\
Language binary Eng. & 5 & 0.842 & 3 & 0.923 & 0.923 \\
Language binary Eng. & 8 & 0.842 & 2 & 0.962 & 0.962 \\
Language binary Eng. & 11 & 0.850 & 21 & 0.660 & 0.660 \\
Continent & 5 & 0.634 & 5 & 0.889 & 0.889 \\
Continent & 8 & 0.633 & 3 & 0.944 & 0.944 \\
Continent & 11 & 0.669 & 32 & 0.651 & 0.651 \\
Country top-10 & 5 & 0.781 & 19 & 0.729 & 0.729 \\
Country top-10 & 8 & 0.777 & 19 & 0.784 & 0.784 \\
Country top-10 & 11 & 0.810 & 73 & 0.458 & 0.458 \\
Language top-10 & 5 & 0.716 & 11 & 0.770 & 0.770 \\
Language top-10 & 8 & 0.728 & 10 & 0.895 & 0.895 \\
Language top-10 & 11 & 0.759 & 67 & 0.547 & 0.547 \\
\bottomrule
\end{tabular}
\end{table}

Layer 11 shows consistently lower rates across all probe types, with substantially more split features identified at deeper layers.

\subsubsection{Shuffled and Random Controls}

Both the shuffled-hierarchy control (randomized city-attribute mappings) and the random probe direction control produce 100\% absorption rates across all layers and probe types ($n = 5$ trials each, standard deviation $= 0$).
The dominance-based metric does not discriminate real from shuffled hierarchies.
This result is the central finding of the cross-domain experiment.

\subsubsection{Cosine-Calibrated Metric}

The cosine-calibrated variant requires that the dominant non-split feature have decoder direction cosine similarity $> \tau_{\text{cos}}$ with the probe direction.
At threshold $\tau_{\text{cos}} = 0.1$, the absorption rate is 0\% across all 9 layer-domain combinations that pass the probe quality gate.
Even at $\tau_{\text{cos}} = 0.05$, only 2 of 9 configurations detect a single absorbed instance each.

The discrepancy between dominance-based (11.3--96.2\%) and cosine-calibrated (0\%) rates reveals that GPT-2 Small's 24k SAE does not encode knowledge-hierarchy probe directions as dedicated latents.
With 98\% dead features on city prompts, the surviving features include polysemantic ``super-absorbers'' that dominate the ablation landscape regardless of probe type.

\subsection{Absorption Scaling Surface: Significant Interaction (H3)}

\subsubsection{Setup}

We analyze 420 SAEs from the SAEBench precomputed dataset (Gemma 2 2B), spanning 9 release families, 3 layers (5, 12, 19), dictionary widths 2,304--1,048,576, and $L_0$ 9.3--8,277.

\subsubsection{Model Comparison}

Three regression models are compared in Table~\ref{tab:scaling}.
The interaction GAM outperforms both simpler models on both $R^2$ and AIC.

\begin{table}[t]
\caption{Scaling surface model comparison. The interaction GAM (bold) achieves the best fit with a highly significant interaction term.}
\label{tab:scaling}
\centering
\small
\begin{tabular}{@{}lccc@{}}
\toprule
Model & $R^2$ & AIC & Interaction $p$ \\
\midrule
Linear & 0.488 & $-1930$ & --- \\
Additive GAM & 0.620 & $-844$ & --- \\
\textbf{Interaction GAM} & $\mathbf{0.693}$ & $\mathbf{-917}$ & $\mathbf{3.1 \times 10^{-15}}$ \\
\bottomrule
\end{tabular}
\end{table}

The interaction term $\text{ti}(\log(\text{width}), \log(L_0))$ is highly significant ($p = 3.1 \times 10^{-15}$), confirming that absorption depends on the joint structure of width and $L_0$.
The interaction GAM explains 69.3\% of variance in absorption rate, a 20.5 percentage-point improvement over the linear model.

Linear model coefficients confirm the directional effects: $\log(\text{width}) = +0.054$ (wider SAEs absorb more), $\log(L_0) = -0.014$ (higher $L_0$ absorbs less), layer $= +0.003$ (deeper layers absorb slightly more).

\subsubsection{Scaling Structure}

Figure~\ref{fig:scaling} shows the absorption surface in $(\log_2(\text{width}), \log_2(L_0))$ space.

\begin{figure}[t]
\centering
\includegraphics[width=\linewidth]{figures/fig5_scaling_surface.pdf}
\caption{Left: raw scatter of 420 SAEs colored by absorption score. Right: GAM-predicted absorption contour surface at layer~12 with labeled phase regions. The high-absorption regime occupies the high-width, low-$L_0$ corner; at $L_0 > 14$, absorption is low regardless of width.}
\label{fig:scaling}
\end{figure}

Per-layer analysis confirms consistent directional effects.
Width correlates positively with absorption ($\rho = +0.35$ to $+0.42$, $p < 10^{-4}$ at all layers) and $L_0$ correlates negatively ($\rho = -0.46$ to $-0.49$, $p < 10^{-8}$ at all layers).
The interaction means these marginal effects are not independent: the width effect is strongest at low $L_0$, and the $L_0$ effect is strongest at high width.

At 1M width, mean absorption reaches 0.703 at layer~12 and 0.585 at layer~19.
At 16k width, mean absorption is 0.034--0.056 across layers, an order of magnitude lower.
Within the 1M stratum, absorption ranges from 0.072 ($L_0 = 114$, layer~5) to 0.896 ($L_0 = 9.3$, layer~12), confirming that $L_0$ drives large variation even at fixed width.

\subsubsection{Phase Boundary Detection}

Gradient analysis of the GAM surface identifies a transition zone at $\log_2(L_0) \in [2.7, 3.8]$, corresponding to $L_0 \approx 6.5$--$14$.
The ridge contains 443 points with gradient magnitude exceeding the 70th percentile threshold (0.69), with maximum gradient magnitude of 0.99.
This gradient ridge marks a region of maximum sensitivity rather than a sharp discontinuity: the GAM surface is continuous, but absorption changes most rapidly in this $L_0$ band.

At $L_0 > 14$, absorption is low regardless of width.
At $L_0 < 7$, absorption increases steeply from 16k to 1M width.
The transition zone represents the regime where practitioners face the sharpest absorption-quality tradeoff.

\subsection{Taxonomy Correction: Artifact Revealed (H4)}

\subsubsection{Setup}

We apply frequency-matched correction to the first-letter absorption taxonomy on GPT-2 Small (\texttt{gpt2-small-res-jb} SAE, layer~8).
Parent features are identified via selectivity heuristic.
For each letter, we compare activation magnitudes in letter-specific contexts against non-letter contexts using a WikiText-103 corpus (7,991 sentences processed).

\subsubsection{Results}

The original taxonomy classified 23/26 letters as Type~II (magnitude ratio $< 0.5$) with a comprehensive absorption rate of 92.3\%.
After frequency-matched correction using Strategy~A (non-letter context comparison), 19 letters change classification: the corrected comprehensive rate drops to 19.2\% (95\% CI [3.8\%, 34.6\%]).

The correction reveals that most parent features identified via selectivity heuristic are genuinely letter-specific: they fire almost exclusively on tokens starting with the target letter.
For letter B, the fire rate ratio (letter-context / non-letter-context) is 70.7:1; for letter D, 124.8:1; for letter X, 8,288:1.
When these features do fire in non-letter contexts, they fire at comparable or higher magnitudes (corrected magnitude ratios: B $= 10.6$, D $= 2.9$, X $= 17.6$).
The original Type~II classification was an artifact of comparing letter-context magnitudes against an inappropriate global baseline.

The Chanin false-negative-based metric provides independent validation: absorption detected in 19/26 letters (73.1\%), with per-letter rates ranging from 0\% (K, U, Z) to 100\% (E, O).
This confirms that genuine absorption occurs even though the Type~II magnitude metric is unreliable.

%----------------------------------------------------------------------
\section{Discussion}
\label{sec:discussion}
%----------------------------------------------------------------------

\subsection{The Causal Status of Absorption}

The suppression effect for sparse probing---where controlling for $L_0$ strengthened the partial correlation from $r = -0.664$ to $r = -0.746$ (Table~\ref{tab:partial_corr}; Figure~\ref{fig:partial_corr})---is the central finding of this paper.
In classical regression analysis, a suppression variable shares variance with the predictor that does not predict the outcome; removing this non-predictive variance reveals a stronger true association \citep{conger1974revised}.
$L_0$ correlates with absorption (higher $L_0$ reduces absorption) but also introduces variance unrelated to quality, and controlling for it unmasks the genuine absorption-quality link.
This result directly reverses the prior concern---articulated by \citet{chanin2024absorption} and in every subsequent absorption study---that the absorption-quality correlation might be an artifact of $L_0$.

The Baron-Kenny mediation analysis provides a second, independent line of evidence (Table~\ref{tab:mediation}; Figure~\ref{fig:mediation}).
For SCR, absorption fully mediates the $L_0$-quality pathway: the direct effect drops to $c' = -0.003$ ($p = 0.71$) when absorption enters the model, and the indirect effect bootstrap 95\% CI of $[0.007, 0.048]$ excludes zero (Sobel $z = 3.62$, $p = 2.9 \times 10^{-4}$).
For RAVEL TPP, absorption mediates 54\% of $L_0$'s effect (Sobel $z = 2.08$, $p = 0.037$).
The proportion mediated exceeds 1.0 for SCR because the direct effect reverses sign once absorption is controlled---a pattern called inconsistent mediation \citep{mackinnon2007mediation} that arises when the suppression variable's non-predictive variance inflates the total effect.

Rosenbaum sensitivity analysis under Mahalanobis matching (17 pairs; Table~\ref{tab:rosenbaum}) corroborates these results.
The $\Gamma$ parameter reaches 2.65 for RAVEL TPP and 1.85 for sparse probing F1.
A $\Gamma$ of 2.65 means an unmeasured confounder would need to increase the odds of being in the high-absorption group by a factor of 2.65 to nullify the observed quality difference, indicating relatively strong robustness to hidden bias \citep{rosenbaum2002observational}.

A systematic assessment using all nine Bradford Hill criteria for epidemiological causation \citep{hill1965environment} yields 3 strong (strength of association, plausibility, coherence), 5 moderate (consistency, specificity, dose-response, experiment, analogy), and 0 weak criteria.

\begin{table}[t]
\caption{Bradford Hill criteria assessment for the absorption-quality causal claim.}
\label{tab:bradford}
\centering
\small
\begin{tabular}{@{}llp{7cm}@{}}
\toprule
Criterion & Assessment & Key Evidence \\
\midrule
Strength & \textbf{Strong} & Partial $r = -0.746$ (sparse probing), $-0.570$ (SCR) after controlling for width, layer, and $L_0$ \\
Consistency & Moderate & 4/4 analytical methods support the link for 3/4 quality metrics \\
Specificity & Moderate & Absorption affects sparse probing and TPP most consistently; unlearning shows no association \\
Temporality & Plausible & $L_0$ is set before training; absorption emerges during training; quality is measured post-hoc \\
Dose-response & Moderate & 1M stratum: $\rho = -0.480$ ($p = 0.044$) for sparse probing; pooled TPP $\rho = -0.524$ \\
Plausibility & \textbf{Strong} & Mechanistic: when a specific latent subsumes a general one, the general latent's false negatives directly degrade sparse probing and TPP \\
Coherence & \textbf{Strong} & Consistent with \citet{chanin2024absorption}'s original finding, SAEBench's inclusion of absorption as a metric \\
Experiment & Moderate & No randomized experiment possible; mediation provides quasi-experimental evidence; $\Gamma = 2.65$ \\
Analogy & Moderate & Analogous to immunodominance in immunology \\
\bottomrule
\end{tabular}
\end{table}

\textbf{The within-width matching null is an important caveat.}
Exact-width matching (4 pairs) and within-width median split matching (23 pairs) both fail to detect significant quality differences between high- and low-absorption SAEs ($\Gamma = 1.0$ for all metrics; Table~\ref{tab:rosenbaum}).
The causal evidence is strongest in cross-width comparisons, where absorption varies jointly with width and $L_0$.
Whether absorption causes quality degradation within a fixed dictionary width remains unresolved with the current sample size ($n = 15$--$18$ per stratum).

\textbf{Unlearning is unrelated to absorption.}
Across all analyses---partial correlation, mediation, stratified analysis, and Rosenbaum matching---unlearning shows no meaningful association with absorption.
This null result demonstrates that the absorption-quality link is specific to certain quality dimensions, not a global artifact.

\subsection{Why the Dominance Metric Fails on Knowledge Hierarchies}

The cross-domain experiment (Section~\ref{sec:crossdomain_results}; Figure~\ref{fig:crossdomain}; Table~\ref{tab:crossdomain}) produced an unexpected result: the dominance-based absorption metric registers 11.3--96.2\% absorption across all knowledge domains and layers, but shuffled controls show 100\% absorption.
A metric that cannot distinguish real from randomized hierarchies does not measure what it claims to measure.

The root cause is specific to the model-SAE combination used.
GPT-2 Small's 24k-feature SAE has 98\% dead features on city prompts.
Among the approximately 500 surviving features per layer, a small number of polysemantic ``super-absorbers'' dominate---features that activate broadly on tokens related to cities, locations, and named entities, regardless of the specific geographic attribute being probed.

The cosine-calibrated metric returns 0\% absorption across all domains, layers, and threshold values ($\tau_{\text{cos}} \in \{0.05, 0.1, 0.15, 0.2, 0.3\}$).
No SAE feature in the GPT-2 Small dictionary aligns with any knowledge-hierarchy probe direction.
Binary country probes achieve 86--93\% accuracy in the residual stream, confirming that the model encodes geographic information, but this information is distributed across many features in a way that no single SAE latent captures.

\textbf{Recommendation for future work}: Cross-domain absorption studies must use models where knowledge features are linearly separable in SAE decoder space, which likely requires models of at least 2B parameters with SAE widths of 65k or more.

\subsection{Practical Implications of the Scaling Surface}

The 420-SAE absorption scaling surface (Section~\ref{sec:scaling}; Figure~\ref{fig:scaling}) provides the first empirical map for practitioners selecting SAE hyperparameters with absorption risk in mind.

\textbf{$L_0 > 14$ as a practical threshold.}
The GAM gradient analysis identifies a transition zone at $L_0 \approx 6.5$--$14$ where absorption rises sharply.
At $L_0 > 14$, absorption remains below 5\% regardless of dictionary width, even at 1M latents.
At $L_0 < 7$, absorption increases from a mean of 3.4\% at width 16k to 70.3\% at width 1M (layer~12).

\textbf{Width scaling requires concurrent $L_0$ scaling.}
The highly significant width--$L_0$ interaction ($p = 3.1 \times 10^{-15}$) means that increasing dictionary width alone---the standard strategy for improving SAE reconstruction quality and feature specificity---increases absorption risk when $L_0$ is not simultaneously adjusted.

\textbf{A three-regime picture of SAE failure modes.}
Combining the absorption scaling surface with the hedging results of \citet{chanin2025hedging} \emph{suggests} a three-regime map:
\begin{enumerate}
\item \textbf{Hedging regime} (low width, any $L_0$): The dictionary lacks capacity; the SAE merges correlated concepts into polysemantic ``hedge'' latents.
\item \textbf{Absorption regime} (high width, low $L_0$): The dictionary has capacity, but the sparsity constraint forces competition among hierarchically related concepts; specific latents subsume general ones.
\item \textbf{Recovery regime} (high width, high $L_0$): The dictionary has both capacity and sufficient active features per input; absorption and hedging are both low.
\end{enumerate}

The optimal operating point lies in the recovery regime.
The transition boundary at $L_0 \approx 7$--$14$ marks the absorption-to-recovery boundary.

\subsection{Limitations}
\label{sec:limitations}

\textbf{Sample size and architecture scope.}
Phase~1 analyses use 48 SAEs from a single model (Gemma 2 2B) and architecture family (JumpReLU and standard SAEs from Gemma Scope and SAEBench).
Cross-architecture validation with TopK, Gated, Matryoshka, or OrtSAE architectures is needed before the absorption-quality causal claim can be considered general.

\textbf{Cross-domain experiment limited to GPT-2 Small.}
The cross-domain absorption measurement was conducted on GPT-2 Small (124M parameters) due to Gemma 2B HuggingFace access restrictions.
GPT-2 Small's 24k SAE with 98\% dead features on city prompts is not representative of the SAE configurations where absorption was originally documented (Gemma Scope 16k--1M on Gemma 2 2B).
The 0\% cosine-calibrated absorption rate may reflect GPT-2 Small's limited factual knowledge capacity rather than a genuine absence of knowledge-hierarchy absorption.

\textbf{Observational design.}
The mediation analysis assumes a specific causal ordering: $L_0$ (training hyperparameter) $\to$ absorption (emergent property) $\to$ quality (downstream evaluation).
While this ordering is plausible, all measurements are cross-sectional.
An interventional study that modifies absorption while holding all other hyperparameters fixed---using, for example, the OrtSAE orthogonality penalty \citep{korznikov2025ortsae} or masked regularization \citep{narayanaswamy2026masked}---would provide stronger causal evidence.

\textbf{Taxonomy correction reveals diagnostic complexity.}
The corrected taxonomy analysis found that 19/26 letters changed classification when proper comparison baselines were established.
The corrected comprehensive rate dropped from 92.3\% to 19.2\%, while the Chanin false-negative-based absorption rate independently validates absorption in 73.1\% of letters.
This discrepancy between magnitude-ratio-based (19.2\%) and false-negative-based (73.1\%) detection highlights that different operational definitions of ``absorption'' can yield dramatically different prevalence estimates.

\textbf{Suppression effect interpretation.}
The sparse probing suppression effect is statistically robust but its interpretation depends on the assumed causal structure.
If $L_0$ affects quality through pathways other than absorption (e.g., $L_0$ affects reconstruction quality, which independently affects sparse probing), then the ``suppression'' could instead reflect an omitted variable bias in the opposite direction.

%----------------------------------------------------------------------
\section{Conclusion}
\label{sec:conclusion}
%----------------------------------------------------------------------

Controlling for $L_0$ strengthens rather than weakens the absorption-quality link, absorption measurement does not transfer to knowledge hierarchies without metric calibration, and width and $L_0$ interact nonlinearly in a 420-SAE scaling surface.

\textbf{Absorption survives the $L_0$ confound---and the link strengthens.}
Across 48 Gemma Scope SAEs, 3 of 4 downstream quality metrics retain meaningful partial correlations with absorption after controlling for $L_0$, width, and layer.
The sparse probing partial correlation strengthened from $r = -0.664$ to $r = -0.746$ when $L_0$ was added (Table~\ref{tab:partial_corr}), revealing a classical suppression effect.
Baron-Kenny mediation analysis establishes absorption as the primary pathway through which $L_0$ affects SCR (full mediation; Table~\ref{tab:mediation}) and TPP (proportion mediated $= 0.54$).
Rosenbaum sensitivity analysis yields $\Gamma = 2.65$ for TPP (Table~\ref{tab:rosenbaum}).
The within-width matching null remains an important caveat.

\textbf{The dominance-based absorption metric does not transfer to knowledge hierarchies.}
Shuffled-hierarchy controls yield 100\% absorption, indistinguishable from real hierarchies.
A cosine-calibrated variant detects 0\% absorption across all thresholds (Section~\ref{sec:crossdomain_results}; Figure~\ref{fig:crossdomain}).

\textbf{The 420-SAE absorption scaling surface reveals significant width--$L_0$ interaction.}
The interaction GAM achieves $R^2 = 0.693$ with a highly significant tensor interaction term ($p = 3.1 \times 10^{-15}$; Table~\ref{tab:scaling}; Figure~\ref{fig:scaling}).
A transition zone at $L_0 \approx 7$--$14$ separates low- and high-absorption regimes.

\textbf{Taxonomy correction reveals an inflated baseline.}
Frequency-matched correction reduces the original 92.3\% comprehensive absorption rate to 19.2\%.
The Chanin false-negative metric independently validates 73.1\% absorption.

\textbf{Practitioner recommendations:}
(1)~Target $L_0 > 14$ to stay in the low-absorption regime.
(2)~When scaling dictionary width, scale $L_0$ concurrently to avoid amplifying absorption.
(3)~Use cosine-calibrated absorption metrics rather than dominance-only metrics when evaluating SAEs on non-spelling tasks.

\textbf{Future work.}
The confound analysis is limited to 48 SAEs from Gemma 2 2B; cross-architecture replication on Llama, Qwen, and Mistral SAEs is needed.
The cross-domain experiment requires repetition on larger models with dedicated knowledge features.
The most pressing follow-up is an interventional experiment: deliberately reducing absorption (e.g., via OrtSAE or Matryoshka training) and measuring whether downstream quality improves.

\bibliographystyle{plainnat}
\bibliography{references}

%%%%%%%%%%%%%%%%%%%%%%%%%%%%%%%%%%%%%%%%%%%%%%%%%%%%%%%%%%%%
\appendix

\section{NeurIPS Paper Checklist}

\begin{enumerate}

\item {\bf Claims}
    \item[] Question: Do the main claims made in the abstract and introduction accurately reflect the paper's contributions and scope?
    \item[] Answer: \answerYes{}
    \item[] Justification: The abstract and introduction state four contributions (confound resolution, cross-domain measurement, scaling surface, taxonomy correction), each with specific quantitative claims that are supported by experimental results in Sections 4.1--4.4.

\item {\bf Limitations}
    \item[] Question: Does the paper discuss the limitations of the work performed by the authors?
    \item[] Answer: \answerYes{}
    \item[] Justification: Section 5.4 discusses five specific limitations: sample size and architecture scope, GPT-2 Small fallback for cross-domain experiments, observational design, taxonomy diagnostic complexity, and suppression effect interpretation.

\item {\bf Theory Assumptions and Proofs}
    \item[] Question: For each theoretical result, does the paper provide the full set of assumptions and a complete (and correct) proof?
    \item[] Answer: \answerNA{}
    \item[] Justification: This paper is empirical; no new theoretical results are claimed. The mediation analysis assumptions (causal ordering) are stated in Section 5.4.

\item {\bf Experimental Result Reproducibility}
    \item[] Question: Does the paper fully disclose all the information needed to reproduce the main experimental results?
    \item[] Answer: \answerYes{}
    \item[] Justification: Section 3.6 lists all software dependencies and versions. Random seeds, threshold values, sample sizes, and SAE identifiers are specified throughout. All analyses use publicly available SAEBench data and open-source models.

\item {\bf Open access to data and code}
    \item[] Question: Does the paper provide open access to the data and code?
    \item[] Answer: \answerNo{}
    \item[] Justification: Code will be released upon acceptance. All input data (SAEBench scores, Gemma Scope SAEs, GPT-2 Small) are already publicly available.

\item {\bf Experimental Setting/Details}
    \item[] Question: Does the paper specify all the training and test details necessary to understand the results?
    \item[] Answer: \answerYes{}
    \item[] Justification: Section 3 specifies all experimental settings including dataset sizes, model specifications, hyperparameters, train/test splits, and evaluation metrics.

\item {\bf Experiment Statistical Significance}
    \item[] Question: Does the paper report error bars suitably and correctly defined?
    \item[] Answer: \answerYes{}
    \item[] Justification: BCa bootstrap 95\% confidence intervals (10,000 resamples) are reported for all partial correlations (Table 1), width-stratified correlations (Table 2), and mediation indirect effects (Table 3). The Rosenbaum $\Gamma$ parameter quantifies sensitivity to hidden confounders (Table 4).

\item {\bf Experiments Compute Resources}
    \item[] Question: Does the paper provide sufficient information on the computer resources?
    \item[] Answer: \answerYes{}
    \item[] Justification: All analyses are training-free and operate on precomputed SAEBench results. The cross-domain experiment uses GPT-2 Small (124M parameters), requiring a single GPU. Total compute is modest ($< 10$ GPU-hours).

\item {\bf Code Of Ethics}
    \item[] Question: Does the research conform with the NeurIPS Code of Ethics?
    \item[] Answer: \answerYes{}
    \item[] Justification: This work analyzes existing public SAE evaluation data and does not involve human subjects, private data, or dual-use concerns.

\item {\bf Broader Impacts}
    \item[] Question: Does the paper discuss both potential positive societal impacts and negative societal impacts?
    \item[] Answer: \answerNA{}
    \item[] Justification: This is a methodological paper about evaluating sparse autoencoders for mechanistic interpretability. The work improves our ability to assess interpretability tools, which has indirect positive societal impact through better AI transparency, but no direct negative societal impacts.

\item {\bf Safeguards}
    \item[] Question: Does the paper describe safeguards for responsible release?
    \item[] Answer: \answerNA{}
    \item[] Justification: No models or datasets with misuse risk are released.

\item {\bf Licenses for existing assets}
    \item[] Question: Are the creators of assets properly credited and licenses respected?
    \item[] Answer: \answerYes{}
    \item[] Justification: All datasets (SAEBench, RAVEL, WikiText-103), models (GPT-2, Gemma 2), and tools (SAELens, TransformerLens) are cited with original papers.

\item {\bf New Assets}
    \item[] Question: Are new assets well documented?
    \item[] Answer: \answerNA{}
    \item[] Justification: No new datasets or models are released.

\item {\bf Crowdsourcing and Research with Human Subjects}
    \item[] Question: Does the paper include details about crowdsourcing experiments?
    \item[] Answer: \answerNA{}
    \item[] Justification: No human subjects research is conducted.

\item {\bf Institutional Review Board (IRB) Approvals}
    \item[] Question: Does the paper describe IRB approvals?
    \item[] Answer: \answerNA{}
    \item[] Justification: No human subjects research is conducted.

\end{enumerate}

\end{document}
